\documentclass[11pt]{article}

\usepackage[margin=1in]{geometry}
\usepackage{amsmath,amssymb,amsthm,mathtools}
\usepackage{enumitem}
\usepackage{microtype}
\usepackage[hidelinks]{hyperref}
\usepackage{xurl}
\usepackage[nameinlink,noabbrev]{cleveref}

\newtheorem{definition}{Definition}
\newtheorem{remark}{Remark}

\newcommand{\Hh}{\mathsf{H}}
\newcommand{\TV}{\mathrm{TV}}

\title{Transparency Primitives for Proof-Carrying Anonymity Claims:\\ Beacons, VRFs, and Append-Only Logs (Series Synthesis)}
\author{}
\date{Unified archive note v0.03 (2026-02-27; archive v121)}

\begin{document}
\maketitle

\begin{abstract}
Several papers in this archive rely on the same publication/anti-equivocation substrate:
(i) a public randomness source (beacon or VRF),
(ii) a precommitment digest for a release/evaluation manifest, and
(iii) an append-only transparency log with inclusion and consistency proofs.
To keep those papers brief and non-redundant, this note states the minimal interface and security properties assumed.
It is deliberately lightweight: we are not proposing new primitives, only an audit-friendly \emph{contract} that existing systems (Certificate Transparency-style logs, Sigstore/Rekor-like services, and beacon/VRF sources) can satisfy.
\end{abstract}

\paragraph{Scope.}
This note specifies the \emph{assumed interface} for public randomness and append-only publication used by the evaluation and release-workflow papers.

\paragraph{Imports.}
We treat cryptographic constructions as black boxes and cite representative standards (VRFs and CT-style logs).
The artifact-policy note covers what is shipped inline vs.\ on-request.\cite{series:artifactpolicy}

\paragraph{Exports.}
We export (i) a minimal beacon/VRF binding contract, (ii) an append-only log contract (entries, checkpoints, inclusion/consistency proofs), and (iii) a small set of failure modes that must be priced in threat-window declarations.

\section{Public randomness: beacons and VRFs}
Many evaluation workflows need randomness that the evaluator cannot grind after seeing outcomes.
We assume a \emph{public randomness interface} that yields values that are (a) unpredictable before a declared time and (b) publicly verifiable after publication.

\begin{definition}[Beacon interface]\label{def:beacon}
A beacon is a public function that maps a round identifier $r$ to an output $R_r\in\{0,1\}^\ell$ together with verification data.
The interface provides:
(i) \textsf{VerifyBeacon}$(r,R_r,\pi_r)\in\{0,1\}$ and
(ii) a public schedule indicating when $R_r$ becomes available.
Security assumptions:
\begin{itemize}[leftmargin=1.2em]
\item \textbf{(Unpredictability)} Before the disclosure time of round $r$, no efficient party can predict $R_r$ with non-negligible advantage.
\item \textbf{(Public verifiability)} Anyone can verify that $(r,R_r)$ is valid.
\end{itemize}
Examples include drand-style beacons and NIST randomness beacons.\cite{drand, nist_beacon}
\end{definition}

\begin{definition}[VRF interface]\label{def:vrf}
A VRF provides a keypair $(\mathsf{pk},\mathsf{sk})$ and algorithms
$\mathsf{Eval}(\mathsf{sk},m)\to (y,\pi)$ and $\mathsf{Verify}(\mathsf{pk},m,y,\pi)\in\{0,1\}$,
where $y$ is pseudorandom conditioned on $\mathsf{pk}$ and $m$.
Security assumptions:
\begin{itemize}[leftmargin=1.2em]
\item \textbf{(Uniqueness)} For a fixed $(\mathsf{pk},m)$ there is at most one valid output $y$ (prevents equivocation).
\item \textbf{(Pseudorandomness)} Without $\mathsf{sk}$, $y$ is indistinguishable from random given $(\mathsf{pk},m)$.
\end{itemize}
See, e.g., the VRF standard.\cite{rfc9381_vrf}
\end{definition}

\section{Binding: manifest digests to evaluation seeds}
The core anti-grinding pattern is: \emph{commit first, then derive randomness from a public source.}

\begin{definition}[Canonical digest]\label{def:digest}
A manifest (release plan, evaluation plan, parameter list) is mapped to a byte string by a canonicalization rule and hashed:
\[
D \;=\; \Hh(\textsf{Canonicalize}(\textsf{manifest})).
\]
The only requirement is that independent verifiers compute the same digest $D$ from the published manifest bytes.\cite{rfc8785_jcs}
\end{definition}

\begin{definition}[Seed binding]\label{def:seedbinding}
Given a manifest digest $D$ and a public randomness value $U$ (beacon output $R_r$ or VRF output $y$), define the evaluation seed
\[
\textsf{Seed} \;=\; \Hh(D\parallel U\parallel \textsf{domain\_sep}),
\]
where \textsf{domain\_sep} prevents cross-protocol reuse.
A valid certificate must disclose $(D,U,\textsf{Seed})$ and the verification data for $U$.
\end{definition}

\begin{remark}[Threat-window implications]
If an evaluator can choose $U$ \emph{after} observing partial outcomes (e.g., select the beacon round post-hoc), grinding reappears.
Threat-window declarations should therefore treat ``choice of $U$'' as a control-plane action that must be precommitted and logged.
\end{remark}

\section{Append-only transparency logs}
A transparency log is a public, append-only data structure that supports efficient proofs of inclusion and append-only consistency.

\begin{definition}[Log interface]\label{def:log}
A log accepts entries $e$ and returns:
(i) an \emph{inclusion proof} $\pi_{\mathrm{inc}}$ that $e$ is included at tree size $n$ under root $h_n$,
and (ii) periodic \emph{checkpoints} (tree size $n$, root $h_n$, timestamp) signed by the log operator.
The interface supports \textsf{VerifyInclusion}$(e,n,h_n,\pi_{\mathrm{inc}})$ and
\textsf{VerifyConsistency}$(n_1,h_{n_1},n_2,h_{n_2},\pi_{\mathrm{con}})$ for append-only growth.
Certificate Transparency-style Merkle-tree logs are the canonical example.\cite{rfc9162_ct}
\end{definition}

\begin{definition}[Anti-equivocation goal]\label{def:antiequiv}
A log is \emph{anti-equivocating} for the purposes of this archive if:
(i) checkpoints are widely witnessed/monitored (or cosigned) so split views are detectably inconsistent, and
(ii) certificates carry enough log evidence (checkpoint + inclusion proof) that verifiers can test consistency against the public view.
This is a deployment/interface goal rather than a new cryptographic definition.\cite{c2sp_tlog_witness,newman2022sigstore}
\end{definition}

\paragraph{What gets logged?}
For this archive, a typical entry is a tuple
\[
e = (\textsf{kind}, D, U, \textsf{metadata}),
\]
where \textsf{kind} distinguishes release manifests from evaluation attempts, and $\textsf{metadata}$ includes the declared spending schedule id, attempt index, \emph{exposure normal form id} (the transcript projection being budgeted), and pointer(s) to on-request artifacts.

\begin{remark}[Surface binding is part of anti-grinding]
A log entry that binds code, parameters, and randomness but omits the \emph{claim surface} still leaves room for semantic drift: the publisher can later present the same run as a claim about a different transcript projection.
For certificate-carrying papers in this archive, the minimal safe rule is therefore: if a budget is claimed on a projection $X$, then the logged statement must hash-bind a stable identifier for that projection (an exposure normal-form id / schema hash).
This is what makes later ``same claim versus new claim'' checks mechanically meaningful.
\end{remark}

\section{Failure modes to price}
The interface above is small, but several failure modes must be priced explicitly:
\begin{itemize}[leftmargin=1.2em]
\item \textbf{Split view / inconsistent checkpoints.} Mitigate with witness policies, monitors, or cosigned checkpoints.
\item \textbf{Post-hoc randomness choice.} Require precommitment of beacon round / VRF message and log it before outcomes.
\item \textbf{Canonicalization ambiguity.} Require a single canonical JSON/text scheme and publish the exact bytes that were hashed.
\end{itemize}
These are not ``paper cuts'': they are exactly the channels by which stop-when-pass workflows re-enter.

\section{Takeaways}
\begin{itemize}[leftmargin=*]
\item \textbf{Commit first, randomize second:} derive evaluation seeds from a manifest digest and a public randomness source.
\item \textbf{Logs are the anti-abort primitive:} publish every initiated attempt (including aborts) to make grinding auditable.
\item \textbf{Specify the interface:} papers stay brief if they can import a shared log/beacon contract rather than re-defining it.
\end{itemize}

\begin{thebibliography}{99}\small

\bibitem{series:artifactpolicy}
Anonymous.
\newblock Artifact policy (source-only archive; artifacts available on request).
\newblock Series note (Synthesis~6), 2026.

\bibitem{rfc9162_ct}
IETF.
\newblock RFC 9162: Certificate Transparency Version 2.0.
\newblock 2021.

\bibitem{rfc9381_vrf}
IETF.
\newblock RFC 9381: Verifiable Random Functions (VRFs).
\newblock 2023.

\bibitem{rfc8785_jcs}
IETF.
\newblock RFC 8785: JSON Canonicalization Scheme (JCS).
\newblock 2020.

\bibitem{drand}
drand project.
\newblock Distributed randomness beacon.
\newblock Public documentation, 2020--2026.

\bibitem{nist_beacon}
NIST.
\newblock NIST Randomness Beacon.
\newblock Public documentation, 2013--2026.

\bibitem{c2sp_tlog_witness}
C2SP.
\newblock Transparency log witness policy / witness ecosystem notes.
\newblock Public drafts, 2023--2026.

\bibitem{newman2022sigstore}
C.~Newman et al.
\newblock Sigstore: Software signing for everybody.
\newblock Public report / paper, 2022.

\end{thebibliography}

\end{document}
